% from ega4-2-fr.tex
\oldpage[IV-2]{48}

Supposons d'abord $n=1$~; il n'y a alors qu'un seul idéal premier minimal
$\mathfrak p$ de $A$, et dire que $M_{\mathfrak p}$ est de longueur $1$
signifie que $M_{\mathfrak p}$ est isomorphe au corps résiduel
$k=A_{\mathfrak p}/\mathfrak pA_{\mathfrak p}$ en tant que
$A_{\mathfrak p}$-module. Par suite $M_{\mathfrak p}$ est annulé par
$\mathfrak pA_{\mathfrak p}$, donc $\mathfrak p$ est l'annulateur de $M$
(Bourbaki, \emph{Alg. comm.}, chap.~II, \S~2, n\textsuperscript{o}~4,
formule~(9)), ce qui entraîne $\mathfrak p=0$ puisque $M$ est supposé fidèle~;
l'anneau $A$ est donc intègre. Cela étant, l'hypothèse $M\neq0$ entraîne
$M/tM\neq0$ par le lemme de Nakayama, et par suite
$\operatorname{long}(M/tM)\geqslant1$, ce qui démontre (3.4.1.3) dans ce cas.
En outre, si $\operatorname{long}(M/tM)=1$, $M$ est nécessairement monogène
(Bourbaki, \emph{Alg. comm.}, chap.~II, \S~3, n\textsuperscript{o}~2, cor.~2
de la prop.~4) donc isomorphe à un quotient $A/\mathfrak b$~; comme il est
fidèle, on a nécessairement $\mathfrak b=0$ et $M$ est isomorphe à $A$~;
comme $\operatorname{long}(A/tA)=1$, $tA$ est nécessairement égal à l'idéal
maximal $\mathfrak m$, et comme $A$ est un anneau local intègre noethérien,
cela prouve que $A$ est un anneau de valuation discrète (Bourbaki,
\emph{Alg. comm.}, chap.~VI, \S~3, n\textsuperscript{o}~6, prop.~9), dont $t$
est l'uniformisante. Réciproquement, si $A$ est un anneau de valuation
discrète, $t$ son uniformisante, $\operatorname{long}(M_{\mathfrak p})=1$ et
si $t$ est $M$-régulier, $M$ est sans torsion, donc isomorphe à un sous-module
de $A$ ($M$ étant de type fini), et par suite isomorphe à $A$ lui-même, d'où
$\operatorname{long}(M/tM)=\operatorname{long}(A/tA)=1$.

Supposons maintenant $n\geqslant2$~; il existe alors une suite exacte
\[
0\longrightarrow M'\longrightarrow M\longrightarrow M''\longrightarrow0
\]
avec $M'\neq0$, $M''\neq0$ et
$\operatorname{Ass}(M)=\operatorname{Ass}(M')\cup\operatorname{Ass}(M'')$~;
en effet, si $\operatorname{Ass}(M)$ n'est pas réduit à un seul élément, cela
résulte de Bourbaki, \emph{Alg. comm.}, chap.~IV, \S~1,
n\textsuperscript{o}~1, prop.~4~; si au contraire $\operatorname{Ass}(M)$ est
réduit à un seul idéal premier, ce dernier est nécessairement l'unique idéal
premier minimal $\mathfrak p$ de $A$~; l'hypothèse entraîne alors
$\operatorname{long}(M_{\mathfrak p})\geqslant2$ et il suffit de prendre pour
$M'$ l'image réciproque d'un sous-module de $M_{\mathfrak p}$ distinct de $0$
et de $M_{\mathfrak p}$. Comme $t$ est $M$-régulier, $t$ n'appartient à aucun
des idéaux premiers de $\operatorname{Ass}(M)$ (Bourbaki,
\emph{Alg. comm.}, chap.~IV, \S~1, n\textsuperscript{o}~1, cor.~2 de la
prop.~2), donc, pour la même raison, $t$ est $M'$-régulier et
$M''$-régulier. Cette dernière propriété entraîne par (3.4.1.4) que la suite
\[
0\longrightarrow M'/tM'\longrightarrow M/tM
\longrightarrow M''/tM''\longrightarrow0
\]
est exacte~; comme par ailleurs il en est de même de la suite
\[
0\longrightarrow M'_{\mathfrak p_i}\longrightarrow M_{\mathfrak p_i}
\longrightarrow M''_{\mathfrak p_i}\longrightarrow0
\]
pour tout $i$, on a
\[
\begin{aligned}
\operatorname{long}(M/tM)
&=\operatorname{long}(M'/tM')+\operatorname{long}(M''/tM''),\\
\operatorname{long}(M_{\mathfrak p_i})
&=\operatorname{long}(M'_{\mathfrak p_i})
 +\operatorname{long}(M''_{\mathfrak p_i})
\end{aligned}
\]
et l'hypothèse de récurrence entraîne donc l'inégalité (3.4.1.3). En outre,
les deux membres ne peuvent être égaux que si les inégalités analogues pour
$M'$ et $M''$ sont aussi des égalités. En vertu de l'hypothèse de récurrence,
cela équivaut à la propriété $\beta$) pour les $\mathfrak p_i$ tels que
$M'_{\mathfrak p_i}\neq0$ ou $M''_{\mathfrak p_i}\neq0$~; mais ces idéaux sont
précisément ceux pour lesquels $M_{\mathfrak p_i}\neq0$. C.Q.F.D.

